\subsubsection{An independent four-dimensional character calculation}

The general shear lemma is Lemma~\ref{lem:shear-descent}.
The following character calculation gives a second proof of the
even-dimensional scalar-kernel step when \(n=4\).  Write \(D=\det V\).

For the next lemma write
\([a,b]=\Sym^a U\otimes\Sym^b U'\), where \(U,U'\)
are the two standard two-dimensional representations of
\(SL_2\times SL_2\).  Only even weights occur below, so all
representations descend through its double cover of \(SO_4\).

\begin{lemma}[The exterior-cube character]
\label{lem:orthogonal-exterior-cube}
Let \(q\) be nondegenerate on the four-dimensional space \(V\),
and put \(W=\Sym^2V\).  Under \(SL_2\times SL_2\to SO(V,q)\),
choose \(V=[1,1]\).  Then
\begin{equation}\label{eq:explicit-so4-exterior-cube}
\begin{split}
 \bigwedge^3 W\simeq{}&[6,0]\oplus[0,6]\oplus[4,4]
       \oplus2[4,2]\oplus2[2,4]\\
       &\oplus[2,2]\oplus2[2,0]\oplus2[0,2].
\end{split}
\end{equation}
In particular \((\bigwedge^3W)^{SO(V,q)}=0\), and
\(\Hom_{O(V,q)}(D,\bigwedge^3W)=0\).
\end{lemma}
\begin{proof}
Put \(A_2=\Sym^2U\) and \(B_2=\Sym^2U'\).
The symmetric Cauchy formula gives
\[
 W=(A_2\otimes B_2)\oplus\C\rho,
                 \qquad \rho=q^{-1},
\]
where \(\rho\) is invariant.  Write \(T=A_2\otimes B_2\).
For any two three-dimensional vector spaces, antisymmetrization in
\((A_2\otimes B_2)^{\otimes r}\), for \(r=2,3\), gives
\begin{align*}
 \bigwedge^2T&=(\Sym^2A_2\otimes\bigwedge^2B_2)
                 \oplus(\bigwedge^2A_2\otimes\Sym^2B_2),\\
 \bigwedge^3T&=(\Sym^3A_2\otimes\bigwedge^3B_2)
       \oplus(\mathbb S_{(2,1)}A_2\otimes\mathbb S_{(2,1)}B_2)
       \oplus(\bigwedge^3A_2\otimes\Sym^3B_2).
\end{align*}
For clarity, the reason for the conjugate partitions in these formulas
is that tensoring a symmetric-group representation with the sign
representation conjugates its partition.  Schur--Weyl decomposition
of the two tensor powers therefore gives exactly the displayed
summands.  This is the usual skew Cauchy formula, not an additional
hypothesis about this particular \(q\).

Here are all the required one-factor calculations:
\begin{equation}\label{eq:spin-one-calculations}
\begin{array}{c|c}
\text{functor on }\Sym^2U&\text{restriction to }SL_2\\ \hline
\Sym^2&\Sym^4U\oplus\C\\
\bigwedge^2&\Sym^2U\\
\bigwedge^3&\C\\
\Sym^3&\Sym^6U\oplus\Sym^2U\\
\mathbb S_{(2,1)}&\Sym^4U\oplus\Sym^2U
\end{array}
\end{equation}
The weights of \(\Sym^2U\) are \(2,0,-2\).
Taking unordered pairs or triples of these weights verifies the
symmetric-power rows; taking distinct weights verifies the exterior
rows.  Complete reducibility identifies the resulting characters
with the stated representations.  Alternatively the final row
follows from
\(A_2\otimes\bigwedge^2A_2=
\mathbb S_{(2,1)}A_2\oplus\bigwedge^3A_2\): the tensor product
\(\Sym^2U\otimes\Sym^2U\) is
\(\Sym^4U\oplus\Sym^2U\oplus\C\), and the exterior cube
removes the scalar summand.

Substitution gives
\begin{align*}
 \bigwedge^2T&=[4,2]\oplus[0,2]\oplus[2,4]\oplus[2,0],\\
 \bigwedge^3T&=[6,0]\oplus[2,0]\oplus[4,4]\oplus[4,2]
       \oplus[2,4]\oplus[2,2]\oplus[0,6]\oplus[0,2].
\end{align*}
Since \(\bigwedge^3W=\bigwedge^3T\oplus
(\bigwedge^2T)\wedge\rho\), these equalities prove
\eqref{eq:explicit-so4-exterior-cube}.  The dimensions are
\(84+36=120=\binom{10}{3}\).  There is no summand \([0,0]\).
Finally the determinant character of \(O(V,q)\), namely its
representation on \(D\), restricts to the trivial character of
\(SO(V,q)\).  Any equivariant image of \(D\) would therefore
be a nonzero \(SO(V,q)\)-invariant vector, which the decomposition
excludes.  This tracks the determinant twist without choosing an
extension of each nontrivial \(SO_4\) summand to \(O_4\).
\end{proof}
